% ============================================================================
% Chapter 4: Basics of Software
% ============================================================================

\chapter{Basics of Software}
\label{ch:software-basics}

\begin{objectives}
  \item Define software and explain why computers need it
  \item Distinguish between system software and application software
  \item Give examples of common system and application software
  \item Understand how hardware and software work together
\end{objectives}

\bytetip[tip]{If hardware is the computer's \textbf{body}, then software is its \textbf{mind}. Without software, a computer is just an expensive box of metal and plastic!}

% ---------------------------------------------------------------
\section{What Is Software?}
\label{sec:what-is-software}
\index{software!definition}

\hl{Software} is a set of instructions (programs) that tells the computer hardware what to do. You can't touch software --- it exists as code stored in the computer's memory.

\begin{theorybox}[Definition: Software]
\textbf{Software}\index{software} is a collection of programs, data, and instructions that tell the computer how to perform specific tasks. Unlike hardware, software has no physical form --- you can't pick it up or hold it.
\end{theorybox}

% ---------------------------------------------------------------
\section{Types of Software}
\label{sec:software-types}
\index{software!types}

Software is divided into two main categories:

% --- Figure 4.1: Software Types Pyramid ---
\begin{figure}[H]
\centering
\begin{tikzpicture}
  % Bottom layer - System Software
  \fill[NavyBlue!85] (-5,0) -- (5,0) -- (3.5,2.5) -- (-3.5,2.5) -- cycle;
  \node[text=white, font=\bfseries\sffamily, align=center] at (0,0.8) 
    {System Software --- The Foundation};
  \node[text=white!80, font=\scriptsize\sffamily] at (0,0.2) 
    {Windows \quad|\quad Linux \quad|\quad macOS \quad|\quad Android};
  
  % Top layer - Application Software
  \fill[ModuleBlue] (-3.5,2.5) -- (3.5,2.5) -- (2,4.5) -- (-2,4.5) -- cycle;
  \node[text=white, font=\bfseries\sffamily, align=center] at (0,3.3) 
    {Application Software};
  \node[text=white!80, font=\scriptsize\sffamily] at (0,2.8) 
    {Word \quad|\quad Paint \quad|\quad Games \quad|\quad Calculator};
  
  % Arrow with label
  \draw[FunYellow, line width=2pt, ->, >=Stealth] (5.5,1.0) -- (5.5,3.5);
  \node[right=0.1cm, font=\scriptsize\sffamily, text=NavyBlue, align=left, text width=3.5cm] 
    at (5.7,2.2) {System software loads first, then apps run on top};
\end{tikzpicture}
\caption{Software types: system software supports application software}
\label{fig:software-pyramid}
\end{figure}

\subsection{System Software}
\index{system software}

\textbf{System software}\index{system software} manages the computer's hardware and provides a platform for other software to run. The most important system software is the \textbf{Operating System} (OS)\index{operating system}.

Examples: \textbf{Windows}, \textbf{macOS}, \textbf{Linux}, \textbf{Android}, \textbf{iOS}

\subsection{Application Software}
\index{application software}

\textbf{Application software}\index{application software} (or ``apps'') are programs you use to perform specific tasks --- like writing a document, browsing the internet, or playing a game.

Examples: \textbf{MS Word}, \textbf{MS Paint}, \textbf{Calculator}, \textbf{Google Chrome}, \textbf{Minecraft}

% ---------------------------------------------------------------
\section{Hardware + Software = A Working Computer!}
\label{sec:hw-plus-sw}

Neither hardware nor software can work alone. They need each other!

% --- Figure 4.2: Hardware + Software = Magic! ---
\begin{figure}[H]
\centering
\begin{tikzpicture}[node distance=0.8cm]
  \node[hwbox, minimum width=3cm, fill=HeaderBg] (hw) 
    {\faMicrochip\ Hardware};
  
  \node[right=0.6cm of hw, font=\Huge\bfseries, text=FunPink] (plus) {+};
  
  \node[hwbox, minimum width=3cm, fill=LabGreen!10, draw=LabGreen, right=0.6cm of plus] (sw) 
    {\faCode\ Software};
  
  \node[right=0.6cm of sw, font=\Huge\bfseries, text=FunYellow!80!black] (equals) {=};
  
  \node[hwbox, minimum width=3.5cm, fill=FunPurple!10, draw=FunPurple, right=0.6cm of equals] (result) 
    {\faStar\ Working Computer!};
\end{tikzpicture}
\caption{Hardware and software must work together}
\label{fig:hw-plus-sw}
\end{figure}

\begin{realworld}
Think of it like a music player. The \textbf{hardware} is the speaker and screen. The \textbf{software} is the music app. Without the app, the speaker stays silent. Without the speaker, the app can't play any sound. They \emph{need} each other!
\end{realworld}

% ---------------------------------------------------------------
\section{Common Software Examples}
\label{sec:common-software}

% --- Table 4.1: Common Software Examples ---
\begin{table}[H]
\centering
\caption{Common software examples and what they do}
\label{tab:software-examples}
\renewcommand{\arraystretch}{1.3}
\begin{tabularx}{\textwidth}{l l X}
\toprule
\rowcolor{HeaderRow}
\textcolor{white}{\textbf{Type}} & 
\textcolor{white}{\textbf{Name}} & 
\textcolor{white}{\textbf{What It Does}} \\
\midrule
\rowcolor{RowLight}
System Software & Windows 11 & Manages the computer \\
\rowcolor{RowDark}
System Software & Android & Runs your smartphone \\
\rowcolor{RowLight}
Application & MS Word & Word processing (writing documents) \\
\rowcolor{RowDark}
Application & MS Paint & Drawing \& image editing \\
\rowcolor{RowLight}
Application & Calculator & Maths operations \\
\rowcolor{RowDark}
Application & Chrome & Browse the internet \\
\bottomrule
\end{tabularx}
\end{table}

\bytetip[fun]{Every time you open an app on your phone, you're running \textbf{application software} on top of \textbf{system software} (Android or iOS). It's like a building: the OS is the foundation, and apps are the rooms you actually live in!}

\begin{funfact}
The first ever software bug was an \emph{actual} bug! In 1947, engineers found a moth stuck inside a computer at Harvard University. They taped it into their logbook and wrote ``first actual case of bug being found.'' That's how the term ``computer bug'' was born!
\end{funfact}

% ---------------------------------------------------------------
% End-of-Chapter
% ---------------------------------------------------------------

\begin{reviewqs}
  \item What is software? How is it different from hardware?
  \item Name the two main types of software.
  \item Give three examples of \textbf{system software}.
  \item Give four examples of \textbf{application software}.
  \item Can hardware work without software? Explain your answer.
  \item What is an operating system, and why is it important?
\end{reviewqs}

\begin{challenge}
Make a two-column table in your notebook. In the left column, list \textbf{10 software programs} you've used. In the right column, classify each one as either ``System Software'' or ``Application Software.'' Bonus: add a third column describing what each program does!
\end{challenge}
